\section{Institutional Interpretation and Empirical Content}\label{sec:institutional}

The theory is intended to capture environments in which downstream interaction generates information that improves an upstream product and that improvement can be reused outside the originating route. The institutional evidence assembled before the theory freeze supports this primitive across settings including appliances, automobiles, agricultural machinery, heavy equipment, and medical devices. The evidence is used here only to discipline the economic interpretation. It does not establish the full causal mechanism or identify the elasticity from downstream rents or sales to information-production effort.

A useful interpretation is a downstream service or field node that observes failures, use conditions, customer needs, maintenance patterns, or other product-relevant information. Producing usable information can require costly diagnosis, documentation, communication, or engineering interaction. Once transmitted upstream, however, the resulting design or quality improvement may benefit units sold through multiple routes. In the language of the model, the originating node bears the cost of $e$ and captures only an own-route return, while $g(e)$ can raise $D_R$ as well as $D_S$.

This interpretation differs from a pure demand-information-sharing problem. Information acquisition and sharing in vertical relationships can affect pricing, disclosure, and channel decisions, as emphasized by \citet{Guo2009} and \citet{GuoIyer2010}. Related work also studies retailer information acquisition under manufacturer encroachment and quality investment \citep{HuEtAl2021}, and information sharing with endogenous channel structure and investment spillovers \citep{HuangChenGuan2020}. Our mechanism instead holds the information-production technology and route structure fixed and asks how an exogenous reallocation of downstream activity changes the incentive to generate a reusable product-improvement input.

\subsection{Frozen empirical predictions}

The model yields five interpretation and falsifiability predictions. They are not an empirical design, and no additional predictions are needed for the theoretical argument.

First, private information effort at the information-producing node may decline as downstream activity reallocates away from it. The relevant outcome is not simply the node's sales; it is costly activity devoted to producing information useful for upstream improvement.

Second, manufacturer or system value placed on reusable information may rise at the same time. Thus a declining local incentive can coexist with increasing upstream demand for diagnostics, feedback, or engineering information.

Third, the divergence should be larger in environments where field information is more reusable across routes or products. In the model, greater cross-route marginal value raises the positive term $A_R\omega_R'D_R'$ in condition~\eqref{eq:G}.

Fourth, upstream investment in diagnostic, reporting, or engineering-feedback infrastructure may rise even while the originating node's sales decline. Such infrastructure is naturally interpreted as an attempt to preserve access to an information input whose system value exceeds the private return captured locally.

Fifth, the sign reversal should weaken or disappear when information has no upstream cross-route use. This is the empirical counterpart of Corollary~\ref{cor:reuse}: if $D_R'=0$, downstream reallocation reduces both private and coordinated information incentives in the separable model.

\subsection{Limits of the bridge}

These predictions deliberately stop short of claiming that any named distribution network implements the model exactly. A field or service channel may perform many functions besides producing information, and downstream reallocation may itself be endogenous in practice. The theory isolates one margin: the location of information production can become misaligned with the location at which information creates value. Establishing the magnitude of that wedge in a particular industry would require separate empirical identification.
